% Part: incompleteness
% Chapter: theories-computability
% Section: first-incompleteness

\documentclass[../../../include/open-logic-section]{subfiles}

\begin{document}

\olfileid{inc}{tcp}{inc}

\olsection{$\Th{Q}$ ora nduweni Ekstensi kang Jangkep, Konsisten, lan
  !!^{axiomatizable}}

\begin{thm}
\ollabel{thm:first-incompleteness}
Ora ana ekstensi $\Th{Q}$ kang jangkep, konsisten, lan
!!{axiomatizable}.
\end{thm}

\begin{proof}
Kita wis ngerti yen ora ana ekstensi $\Th{Q}$ kang konsisten lan
bisa diputusake. Nanging yen $\Th{T}$ jangkep lan !!{axiomatized},
mula teori iki bisa diputusake.
\end{proof}

\begin{explain}
Teorema iki ora adoh saka rumusan asli G\"odel taun 1931 ngenani
Teorema Ora Jangkep Kapisan. Saliyane istilah kang luwih modern,
bedane kang pokok mangkene: G\"odel nggunakake
``$\omega$-konsisten'' tinimbang ``konsisten''; lan dheweke durung bisa
nyebut ``!!{axiomatizable}'' kanthi cakupan umum, amarga gagasan formal
ngenani komputabilitas durung ana. (Model formal komputabilitas
dikembangake sajrone dasawarsa sabanjure, kalebu dening G\"odel, lan
utamane supaya bisa nemtokake jinis teori kang kena marang fenomena
G\"odel.)

Teorema iki nyatakake yen ora bisa nduweni kabeh telu bebarengan:
kejangkepan, kekonsistenan, lan !!{axiomatizability}. Nanging yen
salah sijine ditinggalake, loro liyane bisa diduweni: $\Th{Q}$
konsisten lan bisa diaksiomatisasi kanthi komputabel, nanging ora
jangkep; teori kang ora konsisten iku jangkep lan bisa diaksiomatisasi
kanthi komputabel (upamane, dening $\{ \eq/[0][0] \}$), nanging ora
konsisten; lan himpunan ukara aritmetika kang bener iku jangkep lan
konsisten, nanging ora bisa diaksiomatisasi kanthi komputabel.
\end{explain}
\end{document}
